\section{Problem Statement}
\textcolor{blue}{This section needs editing/review}

In this section, we formalize the challenge of a SQL agent $A$ enforcing a policy $P$ over a database $D$ when subjected to an adaptive red-teamer $R$. 

\subsection{Policy Enforcement on Database Resources}

We define the ground-truth policy as a function $P : q \times D \to \{\texttt{allow}, \texttt{deny}\}$ that determines whether a SQL query $q$ is permissible given the state of database $D$. For a fixed $D$, the policy partitions the space of all possible queries into the denied set $Q^P_{\texttt{deny}}$ and the allowed set $Q^P_{\texttt{allow}}$.

The agent $A_P$ acts as the enforcement layer. We say $A_P$ successfully enforces $P$ if, for any user request $r_t$ and history $c_{t-1}$, the agent's internal query generation $q$ satisfies the following consistency with $P$:
\begin{equation}
\text{If } A_P(c_{t-1}, r_t) \rightarrow q, \text{ then } 
\begin{cases} 
q \in Q^P_{\texttt{allow}} & \implies a_t = D(q) \\
q \in Q^P_{\texttt{deny}} & \implies a_t = \bot
\end{cases}
\end{equation}
An enforcement failure occurs if $A_P$ produces a query $q \in Q^P_{\texttt{deny}}$ that is subsequently executed by $D$, resulting in $a_t \neq \bot$.

\subsection{Adaptive Red Teaming}

Following the interaction model in \S\ref{sec:background}, given an objective $i$, the red-teamer $R$ attempts to generate requests $r_t$ that induce a policy violation. 
% A \textit{conversation} is a sequence $c = \bigl(r_1, a_1, \dots, r_{\ell_c}, a_{\ell_c}\bigr)$.
The goal of $R$ is to generate a set of conversations $C_R$ such that at least one turn $t$ yields $q \in Q^P_{\texttt{deny}}$ while $a_t \neq \bot$.

To evaluate the security--utility tradeoff, we also consider a set of benign conversations $C_B$ consisting of legitimate enterprise requests where every intended query $q$ lies in $Q^P_{\texttt{allow}}$. 
We restrict our evaluation to resource-accessing turns, where $a_t \neq \bot$.
% where $A_P$ attempts to generate a query $q$ to interface with $D$. 
Let $\mathcal{T}^\star_R$ and $\mathcal{T}^\star_B$ denote the sets of resource-accessing turns (represented as $(r_t, a_t)$ pairs) for the red-team and benign datasets, respectively.

\subsection{Evaluation Metrics}

We quantify the agent's behavior using two primary metrics:

\textbf{Policy Violation Rate (PVR):} This measures the frequency of successful exploitations by $R$. We define $\mathrm{PVR}_{\text{turn}}$ as the ratio of resource-accessing turns where the agent allows a prohibited query:
\begin{equation}
\label{eq:pvr-turn}
\mathrm{PVR}_{\text{turn}} = \frac{1}{|\mathcal{T}^\star_R|} \sum_{(r_t, c_{t-1}) \in \mathcal{T}^\star_R} \mathbf{1}[A_P(c_{t-1}, r_t) \in Q^P_{\texttt{deny}}]
\end{equation}

We additionally report two related security metrics:

\begin{itemize}
\item \textbf{Conversation-level Policy Violation Rate}  \  $\mathrm{PVR}_{\text{conv}}$, which credits $R$ with success if \textit{any} turn within a conversation $c \in C_R$ yields a violation, reflecting that a single successful exfiltration compromises the session:

    \begin{equation}
    \label{eq:pvr-conv}
    \begin{aligned}
    \mathrm{PVR}_{\text{conv}} = \frac{1}{|C^\star_\mathcal{R}|} \sum_{c \in C^\star_\mathcal{R}} \mathbf{1} \Bigl[ \, \exists t \leq \ell_c : &\ (q^{(t)}, c_{<t}) \in \mathcal{T}^\star_\mathcal{R} \\
    & \wedge \mathcal{A}(q^{(t)}, c_{<t}, P) = \text{allow} \, \Bigr]
    \end{aligned}
    \end{equation}
    where $C^\star_\mathcal{R} \subseteq C_\mathcal{R}$ is the subset of red-team conversations containing at least one turn in $\mathcal{T}^\star_\mathcal{R}$,

\item \textbf{Work Factor} \ $\mathrm{WF} = \dfrac{1}{\mathrm{PVR}_{\text{turn}}}$, which represents the expected number of resource-accessing attacker turns required to produce one successful policy violation. We report it separately from pVR because PVR (both turn-level and conversation-level variants) compresses towards 0 as defense improves and our adaptive attacker's training budget can be informed by the progress of policy violations (e.g., attacker training should halt when no new policy violations are being discovered).
\end{itemize}

\textbf{Benign Refusal Rate (BRR):} This measures the loss of utility on legitimate requests. $\mathrm{BRR}$ is the rate at which $A_P$ denies ($a_t = \bot$) benign turns. Let $\mathcal{T}_B$ denote the set of all benign turns (resource-accessing or refused) across $C_B$:
\begin{equation}
\label{eq:brr}
\mathrm{BRR} = \frac{1}{|\mathcal{T}_B|} \sum_{(r_t, c_{t-1}) \in \mathcal{T}_B} \mathbf{1}[a_t = \bot]
\end{equation}
We compute BRR over $\mathcal{T}_B$ rather than $\mathcal{T}^\star_B$ so that wrongly-refused benign requests are observable (under $\mathcal{T}^\star_B$ the indicator is identically zero by construction).

Together, these metrics characterize the equilibrium of the deployment. A robust agent must minimize $\mathrm{PVR}$ without a corresponding increase in $\mathrm{BRR}$, ensuring that defensive adaptation does not render the agent unusable for legitimate enterprise tasks.
